\documentclass[10pt,a4paper]{amsart}
\usepackage{amsmath,amssymb,amsthm,graphicx}
% Retained preamble for the Stacks Project chapter snapshot in this directory.
% The chapter's own preamble.tex is a copy of the shared Stacks Project preamble; the block
% below is sliced MECHANICALLY, line for line, from it. Nothing was added, removed or
% redefined; the wrapper class, the xr-hyper/hyperref wiring and the externaldocument list
% are omitted because the assembler supplies the document class.
%
% lines taken from preamble.tex: 3, 5, 9, 11, 12, 156, 157, 158, ... (51 lines total)
%
% package coverage restored from the chapter's own preamble (2026-10-13): xy, xy-2cell

\usepackage[all]{xy}
\xyoption{2cell}
\UseAllTwocells

\newenvironment{reference}{\comment}{\endcomment}
\newenvironment{slogan}{\comment}{\endcomment}
\newenvironment{history}{\comment}{\endcomment}
\theoremstyle{plain}
\newtheorem{theorem}[subsection]{Theorem}
\newtheorem{proposition}[subsection]{Proposition}
\newtheorem{lemma}[subsection]{Lemma}
\theoremstyle{definition}
\newtheorem{definition}[subsection]{Definition}
\newtheorem{example}[subsection]{Example}
\newtheorem{exercise}[subsection]{Exercise}
\newtheorem{situation}[subsection]{Situation}
\theoremstyle{remark}
\newtheorem{remark}[subsection]{Remark}
\newtheorem{remarks}[subsection]{Remarks}
\numberwithin{equation}{subsection}
\def\lim{\mathop{\mathrm{lim}}\nolimits}
\def\colim{\mathop{\mathrm{colim}}\nolimits}
\def\Spec{\mathop{\mathrm{Spec}}}
\def\Hom{\mathop{\mathrm{Hom}}\nolimits}
\def\Ext{\mathop{\mathrm{Ext}}\nolimits}
\def\SheafHom{\mathop{\mathcal{H}\!\mathit{om}}\nolimits}
\def\SheafExt{\mathop{\mathcal{E}\!\mathit{xt}}\nolimits}
\def\Sch{\mathit{Sch}}
\def\Mor{\mathop{\mathrm{Mor}}\nolimits}
\def\Ob{\mathop{\mathrm{Ob}}\nolimits}
\def\Sh{\mathop{\mathit{Sh}}\nolimits}
\def\NL{\mathop{N\!L}\nolimits}
\def\CH{\mathop{\mathrm{CH}}\nolimits}
\def\proetale{{pro\text{-}\acute{e}tale}}
\def\etale{{\acute{e}tale}}
\def\QCoh{\mathit{QCoh}}
\def\Ker{\mathop{\mathrm{Ker}}}
\def\Im{\mathop{\mathrm{Im}}}
\def\Coker{\mathop{\mathrm{Coker}}}
\def\Coim{\mathop{\mathrm{Coim}}}
\DeclareMathSymbol{\boxtimes}{\mathbin}{AMSa}{"02}
\def\QCohstack{\mathcal{QC}\!\mathit{oh}}
\def\Cohstack{\mathcal{C}\!\mathit{oh}}
\def\Spacesstack{\mathcal{S}\!\mathit{paces}}
\def\Quotfunctor{\mathrm{Quot}}
\def\Hilbfunctor{\mathrm{Hilb}}
\def\Curvesstack{\mathcal{C}\!\mathit{urves}}
\def\Polarizedstack{\mathcal{P}\!\mathit{olarized}}
\def\Complexesstack{\mathcal{C}\!\mathit{omplexes}}
\def\Pic{\mathop{\mathrm{Pic}}\nolimits}
\def\Picardstack{\mathcal{P}\!\mathit{ic}}
\def\Picardfunctor{\mathrm{Pic}}
\def\Deformationcategory{\mathcal{D}\!\mathit{ef}}
\usepackage{cleveref}
\crefname{originallemma}{original lemma}{original lemmas}
\Crefname{originallemma}{Original lemma}{Original lemmas}
\makeatletter
\newlabel{descent-lemma-quasi-coherent-alternative-small}{{0GZX}{source}{}{sourcepointer-1}{}}
\newlabel{descent-lemma-quasi-coherent-alternative-small@cref}{{[originallemma][0][]0GZX}{[1][0][]source}{}{}{}}
\makeatother
\crefname{originallemma}{original lemma}{original lemmas}
\Crefname{originallemma}{Original lemma}{Original lemmas}
\makeatletter
\newlabel{equation-alternative-ringed}{{equation-alternative-ringed}{source}{}{sourcepointer-2}{}}
\newlabel{equation-alternative-ringed@cref}{{[originallemma][0][]equation-alternative-ringed}{[1][0][]source}{}{}{}}
\makeatother
\crefname{originallemma}{original lemma}{original lemmas}
\Crefname{originallemma}{Original lemma}{Original lemmas}
\makeatletter
\newlabel{section-restriction}{{section-restriction}{source}{}{sourcepointer-3}{}}
\newlabel{section-restriction@cref}{{[originallemma][0][]section-restriction}{[1][0][]source}{}{}{}}
\makeatother
\crefname{originallemma}{original lemma}{original lemmas}
\Crefname{originallemma}{Original lemma}{Original lemmas}
\makeatletter
\newlabel{topologies-definition-big-small-etale}{{021B}{source}{}{sourcepointer-4}{}}
\newlabel{topologies-definition-big-small-etale@cref}{{[originallemma][0][]021B}{[1][0][]source}{}{}{}}
\makeatother
\title{Stacks, Sheaves on Algebraic Stacks: lemma-locally-quasi-coherent-alternative}
\author{The Stacks Project Authors}
\date{stacks-project:stacks-sheaves.tex:local-snapshot-20261012; GNU FDL 1.2 or later}
\begin{document}
\maketitle
\noindent\textit{Editorial note.} Statement span is lines 4798-4816 of sources/stacks-algebraic-stacks-sheaves/pinned-original-source.tex; proof span is lines 4818-4842, the complete author proof (23 body lines).\par\medskip
\noindent\textit{Editorial note.} Navigation stubs: 4 cited-but-not-included labels are supplied as non-load-bearing pointers with their pinned tags where the snapshot tag index has them.\par\medskip
\noindent\textit{Editorial note.} extra\_packages is empty on purpose: the retained preamble already loads xy with options and a second declaration clashes in the assembler wrapper.\par\medskip
\noindent\textit{Editorial note.} Text is retained verbatim, including the author's own typos and placeholders.\par\medskip
\noindent\textit{Editorial note.} Queue index 19 of queue\_algebraic-stacks-sheaves.json (pb=1296, stmt 4798-4816, proof 4818-4842).\par\medskip
\noindent\textit{Allowed standard prerequisites (editorial).}\begin{itemize}
\item Explicit prerequisite boundary. The cited inputs are prerequisites; the argument in the retained span is not replaced by any cited result.
\end{itemize}\medskip
\ifdefined\hypertarget\hypertarget{sourcepointer-1}{}\fi
\noindent\textit{Source reference.} Original lemma 0GZX: sources/stacks-algebraic-stacks-sheaves/pinned-original-source.tex, source lines 4842--4842. Cited by the retained statement or proof but not part of the collected spans; supplied only as a navigation pointer. The argument does not depend on its content.\par\medskip
\ifdefined\hypertarget\hypertarget{sourcepointer-2}{}\fi
\noindent\textit{Source reference.} Original lemma equation-alternative-ringed: sources/stacks-algebraic-stacks-sheaves/pinned-original-source.tex, source lines 4727--4731. Cited by the retained statement or proof but not part of the collected spans; supplied only as a navigation pointer. The argument does not depend on its content.\par\medskip
\ifdefined\hypertarget\hypertarget{sourcepointer-3}{}\fi
\noindent\textit{Source reference.} Original lemma section-restriction: sources/stacks-algebraic-stacks-sheaves/pinned-original-source.tex, source lines 1185--1200. Cited by the retained statement or proof but not part of the collected spans; supplied only as a navigation pointer. The argument does not depend on its content.\par\medskip
\ifdefined\hypertarget\hypertarget{sourcepointer-4}{}\fi
\noindent\textit{Source reference.} Original lemma 021B: sources/stacks-algebraic-stacks-sheaves/pinned-original-source.tex, source lines 4842--4842. Cited by the retained statement or proof but not part of the collected spans; supplied only as a navigation pointer. The argument does not depend on its content.\par\medskip
\begin{lemma}
\label{lemma-locally-quasi-coherent-alternative}
Let $p : \mathcal{X} \to (\Sch/S)_{fppf}$ be a category fibred in groupoids.
Let $\mathcal{F}$ be an $\mathcal{O}$-module on $\mathcal{X}_{affine}$.
The following are equivalent
\begin{enumerate}
\item for every morphism $x \to x'$ of $\mathcal{X}_{affine}$ such that
$p(x) \to p(x')$ is an \'etale morphism (of affine schemes), the map
$\mathcal{F}(x') \otimes_{\mathcal{O}(x')} \mathcal{O}(x) \to \mathcal{F}(x)$
is an isomorphism,
\item $\mathcal{F}$ is a sheaf for the \'etale topology on
$\mathcal{X}_{affine}$ and for every object $x$ of $\mathcal{X}_{affine}$
the restriction $x^*\mathcal{F}|_{U_{affine, \etale}}$ is quasi-coherent
where $U = p(x)$,
\item $\mathcal{F}$ corresponds to a locally quasi-coherent module
on $\mathcal{X}$ via the equivalence (\ref{equation-alternative-ringed})
for the \'etale topology.
\end{enumerate}
\end{lemma}

\begin{proof}
To make sense out of condition (2), recall that $U_{affine, \etale}$
is the full subcategory of $U_\etale$ consisting of affine objects, see
Topologies, Definition \ref{topologies-definition-big-small-etale}.
Similarly to the discussion in Section \ref{section-restriction}
the object $x$ of $\mathcal{X}_{affine}$ induces an equivalence
$\mathcal{X}_{affine, \etale}/x \to (\textit{Aff}/U)_\etale$ of sites.
Then $x^*\mathcal{F}$ is the sheaf of modules on $(\textit{Aff}/U)_\etale$
corresponding to the restriction
$\mathcal{F}|_{\mathcal{X}_{affine, \etale}/x}$.
Finally, using the continuous and cocontinuous inclusion functor
$U_{affine, \etale} \to (\textit{Aff}/U)_\etale$ we can further
restrict and obtain $x^*\mathcal{F}|_{U_{affine, \etale}}$.

\medskip\noindent
The equivalence of (1) and (2) follows from the remarks above and
Descent, Lemma \ref{descent-lemma-quasi-coherent-alternative-small}
applied to the restriction of $\mathcal{F}$ to $U_{affine, \etale}$
for every object $x$ of $\mathcal{X}$ lying over an affine scheme $U$.
The equivalence of (2) and (3) is immediate from the definitions
and the fact that quasi-coherent modules on $U_{affine, \etale}$
and $U_\etale$ correspond (again by
Descent, Lemma \ref{descent-lemma-quasi-coherent-alternative-small}
for example).
\end{proof}

\end{document}
